\begin{abstract}
Large-scale unsupervised language models (LMs) acquire extensive general knowledge and display some reasoning capabilities; however, precisely directing their output remains a challenge due to the lack of supervision in their training process. Approaches to enhance the controllability of LMs often involve collecting human judgments regarding the relative merit of various model outputs, then adapting the LM through preference alignment, most commonly via reinforcement learning from human feedback (RLHF). Despite their effectiveness, RLHF approaches tend to be intricate and frequently unstable, involving a two-stage process: constructing a reward model that encodes human preferences, then leveraging reinforcement learning to update the LM so as to optimize the induced reward while constraining deviation from the pre-trained distribution. In this work, we propose a novel parameterization for the reward model in RLHF that allows for closed-form derivation of the optimal policy, making it possible to address the standard RLHF problem with a straightforward classification loss. We name our method \textit{Direct Preference Optimization} (DPO), which delivers robust performance with high stability and low computational cost. DPO removes the necessity of sampling from the LM during fine-tuning and dispenses with extensive hyperparameter tuning. Experimental results demonstrate that DPO can adapt LMs to human preferences as effectively as, or better than, prevailing techniques. Remarkably, DPO-based fine-tuning surpasses RLHF methods based on PPO in controlling the sentiment expressed in generated text, and achieves equivalent or superior quality for tasks such as summarization and single-turn dialogue—all while offering significant advantages in simplicity of implementation and training.
\end{abstract}

\section{Introduction}
Large-scale unsupervised language models (LMs) trained on immense text corpora have demonstrated a range of unexpected abilities~\citep{chowdhery2022palm, brown2020language, touvron2023llama, bubeck2023sparks}. Nonetheless, the data these models are exposed to originate from humans who exhibit diverse intentions, preferences, and proficiency levels. This means that certain goals or competencies acquired during training may not be desirable for the model to emulate. For instance, while it is advantageous for an AI programming assistant to \textit{recognize} prevalent coding errors to effectively correct them, we generally want model outputs to be influenced more strongly by examples reflecting the (perhaps infrequent) high-level coding expertise also present in the training data. Likewise, although our language model should be \textit{informed} about a common misconception held by half the population, it should not reiterate that misconception as fact in half of the relevant responses it generates. That is, judiciously \emph{selecting} which \emph{responses and behaviors} a model expresses, out of its broad reservoir of \textit{knowledge and competencies}, is fundamental for building AI systems that are reliable, effective, and governable \citep{ouyang2022training}. While reinforcement learning (RL) is often employed to align LM outputs with human preferences, we demonstrate that the objective conventionally used in RL approaches can actually be achieved precisely by optimizing a basic binary cross-entropy loss, thereby streamlining the process of preference modeling.

Generally, current techniques enforce intended behaviors in a language model using well-constructed collections of human preference data that reflect what is regarded as safe and beneficial by people. This preference learning phase is conducted after the model has undergone large-scale, unsupervised pre-training on diverse textual sources. Although fine-tuning on curated examples of desirable behavior through supervised learning is a natural first approach, reinforcement learning from human (or AI) feedback—RLHF or RLAIF (\citep{christiano2017deep, bai2022constitutional})—has emerged as the leading framework. In RLHF, a reward model is trained on data reflecting human judgments and subsequently guides an LM to generate outputs favored by that reward model, all while ensuring the learned policy remains close to the initial model distribution. RLHF-based strategies yield models with sophisticated conversational and programming performance; however, the RLHF workflow is far more involved than standard supervised training, as it necessitates managing several LMs and performing online sampling and optimization, which introduces considerable computational overhead.

This work introduces a method for optimizing language models to conform to human preferences directly, without requiring explicit reward modeling or reinforcement learning frameworks. We present \textit{{Direct Preference Optimization} (DPO)}, a technique that achieves the same optimization goal as current RLHF approaches—namely, reward maximization subject to a KL-divergence constraint—while maintaining ease of implementation and simplicity in training procedures. Conceptually, DPO works by increasing the log probability ratio of responses that are preferred over those that are not, integrating a dynamically determined, per-example importance weight that mitigates the risk of model collapse often seen with simple probability ratio-based objectives. Similar to other algorithms, DPO employs a formal preference model (e.g., the Bradley-Terry model; \cite{bradley1952rankanalysis}) to evaluate the alignment between a reward function and empirical preference annotations. In contrast to established pipelines—which first utilize this model to train a reward model based on preference losses and then optimize a policy to maximize the reward model—DPO leverages a variable substitution to recast the preference loss directly as a function of the policy parameters. Consequently, given a collection of human-labeled preferences among generated responses, DPO is able to refine the policy using a straightforward binary cross-entropy loss, yielding an optimal policy with respect to a latent reward function tailored to the observed preferences.

Our primary contribution is the introduction of Direct Preference Optimization (DPO), a reinforcement learning-free approach for preference-driven language model training. Experimental results demonstrate that DPO matches or surpasses the efficacy of existing methodologies, including PPO-based RLHF, on a range of tasks such as sentiment control, summarization, and conversational modeling, when applied to language models scaling up to 6B parameters.

Progressively larger self-supervised language models have demonstrated the ability to perform certain tasks without explicit training examples (zero-shot) \citep{radford2019language} or by using only a few provided examples within prompts (few-shot learning) \citep{gpt3,megatron,chowdhery2022palm}. Nevertheless, their effectiveness on downstream applications and their alignment with user objectives are often substantially enhanced through additional fine-tuning, specifically by leveraging corpora containing explicit instructions paired with human-generated responses \citep{mishra-etal-2022-cross,sanh2022multitask,chung2022scaling,thoppilan2022lamda}. This process, referred to as 'instruction-tuning', helps large language models extend their generalization capabilities to new tasks beyond the scope of the training instructions, and typically improves model usefulness and accessibility \citep{chung2022scaling}. While instruction tuning has yielded impressive results, it is typically easier to gather \textit{relative} assessments of response quality from non-experts than to collect full expert-level demonstrations. Building on this, more recent research has leveraged datasets reflecting human preferences to further adjust LLMs, resulting in improved performance in areas such as translation \citep{kreutzer-etal-2018-reliability}, summarization \citep{stiennon2022learning,ziegler2020finetuning}, narrative generation \citep{ziegler2020finetuning}, and instruction-following tasks \citep{ouyang2022training,ramamurthy2023is}. Typically, these methods entail two phases: first, constructing a neural network reward model that captures the preferences embodied in the data, often using a framework such as the Bradley-Terry model \citep{bradley1952rankanalysis}; and second, refining the language model to optimize this learned reward using reinforcement learning algorithms, with approaches such as REINFORCE \citep{williams1992reinforce}, proximal policy optimization (PPO; \cite{schulman2017proximal}), or related methods \citep{ramamurthy2023is}. Closely aligned efforts have made use of LLMs, which have already been tuned for instruction adherence with the aid of human feedback, to generate further synthetic preference datasets—specifically targeting qualities like safety or harmlessness \citep{bai2022constitutional}—by relying on weak supervision from humans, usually limited to rubric-style text guidelines for LLM assessments. These strategies illustrate the intersection of two research trajectories: training language models with reinforcement learning for diverse targets~\citep{Ranzato2015SequenceLT,paulus2018a,wu2018learning}, and the broader area of inferring knowledge from human preference signals \citep{christiano2017deep,kupcsik2018learning}. However, despite the clear advantages of leveraging relative human preferences, the process of applying reinforcement learning to fine-tune large language models continues to pose considerable technical challenges. This paper introduces a theoretically grounded method for directly optimizing relative preference data without the use of reinforcement learning.

Beyond linguistic applications, policy learning from preferences has been explored in both bandit and reinforcement learning paradigms, resulting in a range of methodologies. In particular, contextual bandit problems addressed through preference or ranking feedback, rather than direct reward feedback, fall under the umbrella of contextual dueling bandits (CDB; \cite{yue2012karmed,dudik2015contextual}). When explicit rewards are absent, theoretical treatment of CDBs relies on the concept of a \textit{von Neumann winner}, defined as a policy whose expected success rate against any competing policy is no less than 50\% \citep{dudik2015contextual}. Nevertheless, the CDB paradigm assumes that preference annotations are provided in an online fashion, whereas preference-based policy learning from human annotators often works with a static collection of offline preference-labeled action pairs \citep{yan2022human}. An analogous framework, \textit{preference-based RL} (PbRL), trains policies from binary preference signals produced by an opaque `scoring' function in lieu of explicit rewards \citep{BusaFekete2014,ruiz2023dueling}. A variety of PbRL algorithms have been developed, some of which leverage off-policy preference information, but the majority involve an initial stage where the hidden scoring (reward) function is estimated, followed by an optimization step targeting the learned reward \citep{jain2013learning,BusaFekete2014,christiano2017deep,sadigh2017active,kupcsik2018learning}. In contrast, our approach introduces a unified framework for policy optimization, directly seeking policies that fulfill stated preferences without separately modeling the reward.

\section{Preliminaries}\label{section:prelims}

We begin by summarizing the RLHF pipeline described by \citeauthor{ziegler2020finetuning}, and subsequently adopted in works such as \citep{stiennon2022learning, bai2022training, ouyang2022training}. The standard pipeline comprises three distinct stages: (1) supervised fine-tuning (SFT); (2) preference data collection and reward model construction; and (3) reinforcement learning for policy optimization.

\textbf{SFT}: The RLHF process is typically initiated by adapting a pre-trained language model to the target application domain (e.g., dialogue, summarization) using supervised learning on curated, high-quality examples, yielding the SFT policy $\pi^\text{SFT}$.

\textbf{Reward Modelling Phase}: Subsequently, prompts $x$ are submitted to the SFT-tuned model, resulting in answer pairs $(y_1, y_2)\sim \pi^\text{SFT}(y \mid x)$. Human annotators are then asked to indicate a preference between the two responses, recorded as $y_w\succ y_l \mid x$, where $y_w$ and $y_l$ refer to the favored and disfavored completions, respectively. Although the human judgments are presumed to arise from an underlying, unobservable reward function $r^*(y, x)$, the explicit form of this function remains unknown. A variety of models have been proposed to capture preference data, with the Bradley-Terry (BT) model \cite{bradley1952rankanalysis} being commonly employed (more generally, the Plackett-Luce model \citep{plackett1975analysis, luce2012individual} is applicable for multiple ranked answers). The BT model formulates the human preference probability $p^*$ as follows:

Given a fixed dataset of pairwise preferences $\mathcal{D}=\bigl\{x^{(i)}, y_w^{(i)}, y_l^{(i)}\bigr\}_{i=1}^N$ drawn according to $p^*$, one can define a parameterized reward function $r_{\phi}(x, y)$ and fit its parameters by maximizing the likelihood. This setup can be viewed as a binary classification task, where the objective is to minimize the negative log-likelihood loss:

\begin{equation}\label{eq:reward_model}
    \mathcal{L}_R(r_{\phi}, \mathcal{D}) = -\mathbb{E}_{(x, y_w, y_l)\sim \mathcal{D}}\bigl[\log \sigma(r_{\phi}(x, y_w)- r_{\phi}(x, y_l))\bigr]
\end{equation}

with $\sigma$ representing the logistic function. In large language model applications, $r_{\phi}(x, y)$ is commonly initialized with the SFT model $\pi^\text{SFT}(y \mid x)$, and extended by adding a linear transformation to the top of the model, generating a scalar reward estimate \cite{ziegler2020finetuning}. Previous work also implements a normalization step on the reward outputs to achieve zero mean, ensuring $\mathbb{E}_{x,y\sim \mathcal{D}}\left[r_\phi(x, y)\right] = 0$ over all $x$, as a means to reduce reward variance.

\textbf{RL Fine-Tuning Phase}: During reinforcement learning fine-tuning, this trained reward function provides evaluative feedback to the language model. Consistent with prior studies~\citep{jaques2017sequence, jaques2020human}, the objective is formulated as

\begin{equation}\label{eq:RL}
\max_{\pi_{\theta}}  \mathbb{E}_{x\sim \mathcal{D}, y\sim \pi_{\theta}(y \mid x)}\bigl[r_{\phi}(x, y)\bigr] - \beta\mathbb{D}_{\textrm{KL}}\bigl[\pi_{\theta}(y\mid x)\mid \mid \pi_\text{ref}(y\mid x)\bigr],
\end{equation}

where $\beta$ regulates the penalty for divergence from the reference policy $\pi_\text{ref}$, typically chosen as the initial SFT model $\pi^\text{SFT}$. For practical purposes, $\pi_\theta$ is also initialized from $\pi^\text{SFT}$. The KL-regularization term plays a vital role in preventing the policy from straying excessively from the data regime where the reward model is most reliable, as well as maintaining diversity in generated outputs to avoid collapse onto a few high-reward sequences. Because generating text is a discrete process, direct gradient-based optimization of this objective is infeasible, so reinforcement learning methods are used instead. The established technique~\citep{ziegler2020finetuning, stiennon2022learning, bai2022training, ouyang2022training} introduces a combined reward

${r(x, y) = r_{\phi}(x, y) -\beta (\log \pi_{\theta}(y\mid x) - \log \pi_\text{ref}(y\mid x))}$,

which is maximized via PPO~\cite{schulman2017proximal}.

\section{Direct Preference Optimization}\label{sec:DPO}

Given the practical difficulties associated with scaling reinforcement learning algorithms to large models, such as those used for language generation, our aim is to propose a direct and efficient preference-based policy optimization method. Whereas conventional RLHF approaches require a separate reward modeling phase followed by RL-driven optimization, our technique relies on an alternative reward model parameterization that yields a closed-form solution for its associated optimal policy. In what follows, we elaborate on this approach: by utilizing a well-defined analytic mapping from reward functions to their corresponding optimal policies, we can re-express a reward-model loss in terms of policy parameters. This change-of-variable technique eliminates the need to explicitly train a reward model, yet continues to operate under commonly adopted human preference models, like the Bradley-Terry framework. Consequently, the policy network now serves both as a language generator and as the implicit reward estimator.

\textbf{DPO Objective Derivation.} We begin by considering the standard RL objective given in Eq.~\ref{eq:RL}, utilizing a general reward function $r$. As established in earlier studies~\citep{peters2007reinforcement, peng2019advantage, korbak2022reinforcement, go2023aligning}, it can be shown directly that the policy optimizing the KL-regularized reward maximization problem in Eq.~\ref{eq:RL} can be written as:

\begin{equation}\label{eq:op_policy}
    \pi_r(y\mid x) = \frac{1}{Z(x)}\pi_\text{ref}(y\mid x)\exp\left(\frac{1}{\beta}r(x, y)\right),
\end{equation}

with the normalization constant (partition function) given by $Z(x) =\sum_{y}\pi_\text{ref}(y\mid x)\exp\left(\frac{1}{\beta}r(x, y)\right)$. A detailed derivation is provided in Appendix \ref{app:derivation1}. Even when the reward is approximated via its MLE estimate $r_{\phi}$ for the true reward $r^*$, calculating $Z(x)$ remains computationally burdensome~\citep{korbak2022reinforcement, go2023aligning}, thereby making this form challenging for practical use. However, by manipulating Eq.~\ref{eq:op_policy}, we can instead express the reward function in terms of the optimal policy $\pi_r$, the reference policy $\pi_\text{ref}$, and the unknown $Z(x)$. Specifically, taking the logarithm of Eq.~\ref{eq:op_policy} and rearranging yields:

\begin{equation}\label{eq:main_eq}
    r(x,y) =\beta \log \frac{\pi_r(y\mid x)}{\pi_\text{ref}(y\mid x)} + \beta \log Z(x).
\end{equation}

This alternate form can also be applied to the ground-truth reward $r^*$ and its corresponding optimal policy $\pi^*$. Notably, the Bradley-Terry model considers only the reward difference between two responses; specifically, ${p^*(y_1 \succ y_2 \mid x) = \sigma(r^*(x, y_1) - r^*(x, y_2))}$. If we substitute the reparameterization from Eq.~\ref{eq:main_eq} for $r^*(x,y)$ into the Bradley-Terry preference equation (Eq.~\ref{eq:bradley-terry}), the partition function $Z(x)$ cancels out. As a result, the human preference probability is determined only by the optimal policy $\pi^*$ and the reference $\pi_\text{ref}$. Accordingly, under the Bradley-Terry model, the policy $\pi^*$ solving RLHF must satisfy:

\begin{equation}\label{eq:objective}
    p^*(y_1\succ y_2 \mid x)=\frac{1}{1 + \exp\left(\beta \log \frac{\pi^*(y_2\mid x)}{\pi_\text{ref}(y_2\mid x)} - \beta \log \frac{\pi^*(y_1\mid x)}{\pi_\text{ref}(y_1\mid x)}\right)}
\end{equation}

For a full step-by-step derivation, see Appendix~\ref{app:derivation2}. While Eq.~\ref{eq:objective} is specific to the Bradley-Terry model, this approach extends to Plackett-Luce models~\citep{plackett1975analysis, luce2012individual} as discussed in Appendix~\ref{app:plackett_luce_models}.

Given that human preference probabilities can now be written entirely as functions of the optimal policy, we can formulate a maximum likelihood training criterion directly for a parameterized policy $\pi_\theta$. Drawing an analogy to reward modeling (see Eq.~\ref{eq:reward_model}), the policy-level objective we obtain is:

\begin{equation}\label{eq:optimum_model}
    \mathcal{L}_\text{DPO}(\pi_{\theta}; \pi_\text{ref}) = -\mathbb{E}_{(x, y_w, y_l)\sim \mathcal{D}}\left[\log \sigma \left(\beta \log \frac{\pi_{\theta}(y_w\mid x)}{\pi_\text{ref}(y_w\mid x)} - \beta \log \frac{\pi_{\theta}(

In this formulation, we model an implicit reward through an alternative parameterization, wherein the optimal policy directly corresponds to $\pi_\theta$. Notably, our method aligns with learning a reparametrized Bradley-Terry model, thereby inheriting favorable theoretical guarantees, including consistency given appropriate assumptions on the preference data distribution \cite{bong2022generalized}. Additional theoretical analysis of DPO and its connections to related methods is presented in Section~\ref{sec:theory}.

\textbf{Mechanics of the DPO Update.} To better understand DPO’s inner workings, we can examine the gradient of its loss function, $\mathcal{L}_\text{DPO}$. Taking the gradient with respect to $\theta$ yields:

\begin{multline*}\label{eq:gradient}
    \nabla_\theta \mathcal{L}_\text{DPO}(\pi_\theta;\pi_\text{ref}) = \\ -\beta\mathbb{E}_{(x, y_w, y_l) \sim \mathcal{D}} \bigg[\underbrace{\sigma(\hat{r}_\theta(x, y_l) - \hat{r}_\theta (x, y_w))}_\text{places greater emphasis when reward ranking is inaccurate}\bigg[\underbrace{\nabla_\theta\log \pi(y_w \mid x)}_\text{promotes $y_w$ occurrence} - \underbrace{\nabla_\theta\log\pi(y_l \mid x)}_\text{discourages $y_l$ occurrence}\bigg]\bigg],
\end{multline*}

Here, $\hat{r}_\theta(x, y) = \beta \log \frac{\pi_\theta(y \mid x)}{\pi_\text{ref}(y \mid x)}$ specifies the reward function that the language model $\pi_\theta$ defines in reference to $\pi_\text{ref}$ (with additional discussion in Section~\ref{sec:theory}). Conceptually, the loss gradient for $\mathcal{L}_\text{DPO}$ acts to boost the probability assigned to preferred responses $y_w$ while suppressing that of less desirable responses $y_l$. A crucial aspect of the update is the weighting of examples by the extent to which the implicit reward model $\hat{r}_\theta$ mistakenly scores dispreferred completions more favorably—modulated by $\beta$—reflecting the misordering between candidates while also incorporating the KL constraint’s impact. Empirical findings underscore the significance of this weighting term; eliminating it, as in a basic implementation, may result in pathological model behavior (see Appendix Table~\ref{tab:unlikelihood_generations}).

\textbf{DPO summary.} The standard DPO workflow proceeds as follows: 1) For each input prompt $x$, draw response samples $y_1, y_2 \sim \pi_\text{ref}(\cdot \mid x)$, and then annotate them based on human preference to assemble the offline preference dataset $\mathcal{D} = \{x^{(i)}, y_w^{(i)}, y_l^{(i)}\}_{i=1}^N$; 2) Subsequently, update the parameters of the language model $\pi_\theta$ by minimizing $\mathcal{L}_\text{DPO}$, given $\pi_\text{ref}$, the dataset $\mathcal{D}$, and the specified $\beta$. In real-world settings, there is typically a desire to utilize existing publicly available preference datasets, rather than producing new response samples and collecting fresh preference data. Because such preference datasets are typically constructed using $\pi^\text{SFT}$, we set $\pi_\text{ref} = \pi^\text{SFT}$ whenever this policy is accessible. In scenarios where $\pi^\text{SFT}$ is not at hand, we estimate $\pi_\text{ref}$ by maximizing the likelihood over the set of preferred completions ${(x, y_w)}$, i.e., by setting ${\pi_\text{ref} = \argmax_{\pi}\mathbb{E}_{x, y_w \sim \mathcal{D}}\left[\log \pi(y_w \mid x)\right]}$. Employing this strategy helps to alleviate the distribution mismatch between the unknown true reference policy and the $\pi_\text{ref}$ employed in DPO. Implementation specifics and chosen hyperparameters are provided in Appendix~\ref{app:implementation}.

\section{Theoretical Analysis of DPO} This section presents additional insight into the DPO approach, establishes theoretical justification, and illustrates how DPO addresses challenges inherent in actor-critic methods typically utilized for RLHF (including PPO~\cite{schulman2017proximal}).

\label{sec:theory}

\subsection{Your Language Model Is Secretly a Reward Model} DPO eliminates the necessity of constructing an explicit reward function and executing RL for policy optimization by leveraging a single maximum likelihood formulation. The optimization target described in Eq. \ref{eq:main_eq} aligns with the Bradley-Terry model under a reward parametrization $r^*(x, y) = \beta \log\frac{\pi^*_\theta(y \mid x)}{\pi_\text{ref}(y \mid x)}$, and learning $\pi_\theta$ is equivalent to fitting a reward model as per Eq. \ref{eq:reward_model} after applying a change of variables. In what follows, we develop the underlying theoretical foundation for this reparameterization, demonstrate that it does not impose additional constraints on the expressive power of the learned reward model, and argue that it enables exact retrieval of the optimal policy. We start by introducing a notion of equivalence among reward functions.

\begin{definition}
Two reward functions $r(x, y)$ and $r'(x, y)$ are said to be equivalent if and only if there exists a function $f$ such that ${r(x, y) - r'(x, y) = f(x)}$.
\end{definition}

It is straightforward to verify that this constitutes an equivalence relation, thereby segmenting the space of reward functions into distinct equivalence classes. The following lemmas can thus be stated:

\begin{lemma}\label{lemma:same_prefrence} Within the Plackett-Luce preference model, as well as the Bradley-Terry special case, any two reward functions belonging to the same equivalence class result in identical preference distributions.
\end{lemma}

\begin{lemma}\label{lemma:same_policy}
    Any two reward functions from a single equivalence class will induce the same optimal policy in the context of the constrained RL formulation.
\end{lemma}

The proofs for these claims are uncomplicated, and we provide them in Appendix \ref{app:lemma1}. The first lemma captures a widely recognized issue of under-specification inherent to the Plackett-Luce family \cite{plackett1975analysis}, whereby further identifiability constraints are often necessary in order to obtain meaningful MLE estimates from Eq. \ref{eq:reward_model} \cite{bong2022generalized}. According to the second lemma, the particular choice of representative within an equivalence class does not affect the optimal policy it produces; thus, for our learning objective, any member of the appropriate equivalence class is sufficient. We formalize this in the following theorem, with the proof appearing in Appendix~\ref{app:thm1}:

\begin{theorem}\label{thm:main}
    Provided mild regularity assumptions, all equivalence classes of reward functions that align with the Plackett-Luce (including Bradley-Terry) models can be expressed as ${r(x, y) = \beta \log \frac{\pi(y\mid x)}{\pi_\text{ref}(y\mid x)}}$, where $\pi(y\mid x)$ is some probabilistic model, and $\pi_\text{ref}(y \mid x)$ serves as a fixed reference model.
\end{theorem}

\begin{sproof}
    Let us consider an arbitrary reward function $r(x, y)$ which specifies an optimal policy $\pi_r(y \mid x)$, as given by Eq. \ref{eq:op_policy}. Our aim is to demonstrate that a function belonging to the equivalence class of $r$ admits the proposed reparameterization. To do this, we define the projection
\begin{equation}
    f(r; \pi_\text{ref}, \beta)(x, y) = r(x, y) - \beta\log\sum_{y}\pi_\text{ref}(y\mid x)\exp\left(\frac{1}{\beta}r(x, y)\right)
\end{equation}
This operator $f$ renormalizes the reward function via the log-partition function with respect to the reference policy. Because the normalization depends only on $x$, the result $f(r; \pi_\text{ref}, \beta)(x, y)$ belongs to the equivalence class associated with $r(x, y)$. Substituting $r$ using the right-hand side of Eq.~\ref{eq:main_eq}, valid for any reward function, leads to $f(r; \pi_\text{ref}, \beta)(x, y) = \beta \log \frac{\pi_r(y\mid x)}{\pi_\text{ref}(y\mid x)}$. Thus, the projection $f$ maps $r$ to an equivalent reward function of the stated form, establishing that the reparameterization captures all possible reward classes without loss of generality.
\end{sproof}

An alternative interpretation of Theorem~\ref{thm:main} is that it precisely determines, for each equivalence class, which reward function is chosen by the DPO reparameterization. Specifically, it selects the reward function satisfying:

\begin{equation}\label{eq:lag_p}
     \sum_{y}\underbrace{\pi_\text{ref}(y\mid x)\exp\left(\frac{1}{\beta}r(x, y)\right)}_{=\pi(y\mid x)\text{, using Thm.~\ref{thm:main} reparam.}} = 1,
\end{equation}

meaning that $\pi(y\mid x)$ defines a legitimate probability distribution, since all probabilities are non-negative and sum to one. By referencing Eq.~\ref{eq:op_policy}, we recognize that Eq.~\ref{eq:lag_p} represents the partition function of the optimal policy derived from the reward $r(x, y)$. The DPO algorithm’s main contribution is to introduce specific constraints within the broadly defined Plackett-Luce (and particularly Bradley-Terry) preference models. These constraints do not reduce the expressivity of the reward models but do guarantee that the optimal policy from Eq.~\ref{eq:op_policy} remains analytically manageable across any prompt $x$.

\subsection{Instability of Actor-Critic Algorithms} Our framework also sheds light on potential instability in standard actor-critic methods widely used in RLHF, such as PPO. In this analysis, we adhere to the RLHF methodology and, in particular, focus on the RL fine-tuning phase as detailed in Section \ref{section:prelims}. There are conceptual parallels to be drawn with the control as inference paradigm \cite{levine2018reinforcement} applied to the constrained RL formulation described in \ref{eq:RL}. Here, we consider a parameterized model $\pi_{\theta}(y\mid x)$ and attempt to minimize the divergence $\mathbb{D}_{\text{KL}}[\pi_{\theta}(y|x) \mid \mid \pi^*(y\mid x)]$, with $\pi^*$ representing the optimal policy given by Eq.~\ref{eq:optimum_model}, defined in terms of the reward $r_{\phi}(y, x)$. After appropriate mathematical manipulation, this yields the following optimization objective:

\begin{equation}\label{eq:AC}
    \max_{\pi_{\theta}}\mathbb{E}_{\pi_{\theta}(y\mid x)}\bigg[\underbrace{r_{\phi}(x, y) -\beta\log\sum_{y}\pi_\text{ref}(y\mid x)\exp\left(\frac{1}{\beta}r_{\phi}(x, y)\right)}_{f(r_{\phi}, \pi_\text{ref}, \beta)} - \underbrace{\beta\log\frac{\pi_{\theta}(y\mid x)}{\pi_\text{ref}(y\mid x)}}_{\text{KL}}\bigg]
\end{equation}

The objective under consideration here aligns with those targeted in previous research 
\citep{ziegler2020finetuning, stiennon2022learning, bai2022training, ouyang2022training}, specifically utilizing a DPO-equivalent reward within the $r_{\phi}$ reward class. In this framework, the normalization component of $f(r_{\phi}, \pi_\text{ref}, \beta)$ can be interpreted as representing the soft value function corresponding to the reference policy $\pi_\text{ref}$. Although this normalization does not alter the optimal policy, its absence can lead to high-variance gradients during policy updates, potentially destabilizing training. Incorporating a learned value function is one way to handle normalization, but this approach is itself prone to optimization challenges. Another strategy, observed in prior literature, is to normalize the reward based on a baseline established by human completions, which effectively approximates the normalization factor using a single-sample Monte Carlo method. In contrast, the DPO reparameterization approach directly yields a reward structure that obviates the need for such baselines. 

\section{Experiments}
We now turn to an empirical investigation of DPO's capability to learn policies from preference data. Initially, in a controlled text-generation scenario, we examine DPO’s efficiency in balancing reward maximization against minimizing KL-divergence relative to the reference policy, and compare it to widely-used preference learning approaches like PPO. Subsequently, we assess DPO when applied to more substantial models and challenging RLHF problems, such as tasks in summarization and dialogue systems. Our experiments reveal that, even with minimal hyperparameter adjustment, DPO often matches or surpasses strong benchmarks, including RLHF with PPO and methods selecting the best of $N$ sampled rollouts from a learned reward model. Prior to discussing the outcomes, we provide an overview of our experimental design; further specifics can be found in Appendix~\ref{app:exp_details}.

\textbf{Tasks.} We examine three distinct open-ended text generation tasks in our experimental setup. Across all tasks, the algorithms are trained to optimize a policy using a dataset of preference pairs $\mathcal{D}=\bigl\{x^{(i)}, y_w^{(i)}, y_l^{(i)}\bigr\}_{i=1}^N$. For the \textbf{controlled sentiment generation} task, $x$ represents the beginning of a movie review sampled from the IMDb dataset \cite{maas-EtAl:2011:ACL-HLT2011}, and the policy is required to generate a continuation $y$ exhibiting positive sentiment. To facilitate a controlled assessment, we create preference pairs automatically for generated outputs using a pre-trained sentiment classifier, ensuring that $p(\text{positive}\mid x,y_w)>p(\text{positive}\mid x,y_l)$. For SFT, we perform fine-tuning of GPT-2-large on positive reviews drawn from the IMDb dataset’s training portion, running training until convergence (see further details in App~\ref{app:sentiment_details}). In the \textbf{summarization} task, $x$ consists of a Reddit forum post, and the policy must generate a summary $y$ that distills its key points. We leverage the Reddit TL;DR dataset \citep{volske-etal-2017-tl} for this purpose, utilizing human preference labels collected by \citeauthor{stiennon2022learning}. The corresponding SFT model is fine-tuned on human-created summaries of forum posts,\footnote{\url{https://huggingface.co/CarperAI/openai_summarize_tldr_sft}} and RLHF is performed using the TRLX library \citep{leandro_von_werra_2023_7790115}. The set of human preferences used for evaluation was also gathered by \citeauthor{stiennon2022learning}, though it was based on outputs from a different, but comparably-trained, SFT model. The \textbf{single-turn dialogue} task involves $x$ as a user’s prompt, which could range from technical queries in astrophysics to personal advice. The model is tasked with producing a response $y$ that is informative and engaging. Here, we use the Anthropic Helpful and Harmless dialogue corpus \citep{bai2022training}, which consists of 170,000 dialogues between a user and an automated assistant, each finishing with two model-generated responses and an accompanying human preference annotation. Because no pre-trained SFT model is provided in this context, we construct the SFT model by fine-tuning a generic language model solely on the preferred completions from the dataset.

\textbf{Evaluation.} We employ two primary strategies for assessing our models. To systematically measure each algorithm's capability to solve the constrained reward maximization problem, we operate within a controlled sentiment generation context, where the set of reward-KL-divergence trade-offs attained by each algorithm is calculated; this analysis is enabled by direct access to the true reward signal via a sentiment classifier. Conversely, when shifting to real-world scenarios—where an exact reward function is unavailable—evaluation relies on the \textit{win rate} metric relative to a baseline policy, leveraging GPT-4 as an automated proxy for human assessment. In the summarization domain, the test set’s reference summaries act as the baseline; in single-turn dialogue, the baseline is established by the reference responses favored in the dataset. Although recent findings suggest that large language models can serve as superior automatic evaluators compared to traditional metrics \citep{Chen2023ExploringTU}, we validate the appropriateness of GPT-4 as an evaluator by carrying out a human study, detailed in Sec.~\ref{sec:human-judgments}. Our results indicate a strong alignment between GPT-4 judgments and human preferences, with rates of concordance that are equal to or greater than inter-annotator agreement among human raters.

\textbf{Methods.} Beyond DPO, our evaluation encompasses a suite of established methodologies for guiding language models according to human preferences. For baseline comparisons, we use zero-shot prompting of \textbf{GPT-J} \citep{gpt-j} in summarization and a 2-shot prompting configuration for \textbf{Pythia-2.8B} \citep{biderman2023pythia} in dialogue generation. We further assess the performance of the \textbf{SFT} model and the \textbf{Preferred-FT} variant, which is obtained by supervised fine-tuning on the selected output $y_w$, sourced either from the SFT model (in both sentiment and summarization settings) or from a generic language model (in dialogue tasks). Another method under comparison is \textbf{Unlikelihood}~\citep{welleck2019neural}, which updates the policy by increasing likelihood on $y_w$ while reducing it on $y_l$, optionally weighted by a coefficient $\alpha\in[0,1]$ in the 'unlikelihood' objective term. Additionally, we evaluate \textbf{PPO} \citep{schulman2017proximal} where the reward is learned from preference annotations, as well as \textbf{PPO-GT}, an oracle variant trained directly with the true reward function when available (as in controlled sentiment experiments). For PPO-GT, we use two variants: an off-the-shelf implementation \cite{leandro_von_werra_2023_7790115} and a refined version incorporating reward normalization and tuned hyperparameters for enhanced results (these adjustments are also used for standard PPO with learned rewards). Finally, the \textbf{Best of $N$} baseline samples $N$ candidates from the SFT model (or Preferred-FT in dialogue), returning the one that scores highest according to the preference-trained reward model. Although this approach separates the influence of reward modeling from PPO-style optimization and often yields strong performance, it incurs a significant computational burden, requiring $N$ completions to be generated and scored for each input during testing, making it infeasible for larger $N$.

\subsection{How well can DPO optimize the RLHF objective?}

In standard RLHF approaches, the objective function consists of maximizing rewards while imposing a KL penalty to prevent the policy from straying excessively from the reference policy. Therefore, comparative evaluations between methods must jointly consider both attained rewards and the corresponding KL divergence; a marginally improved reward that accompanies a substantially increased KL value is often not preferable. In Figure~\ref{fig:frontier-tldr-main}, we present the reward versus KL trade-off curves (frontiers) across several algorithms within the sentiment classification framework. For each algorithm, we conduct multiple independent training sessions, where each run uses a unique hyperparameter to control the conservativeness of the policy (for PPO, target KL $\in\{3,6,9,12\}$; for unlikelihood, $\beta \in \{0.05,0.1,1,5\}$ and $\alpha\in\{0.05,0.1,0.5,1\}$; for preferred-FT, distinct random seeds). This experimental design comprises a total of 22 runs. At intervals of 100 training steps, up to convergence, we assess each resulting policy on a collection of held-out test prompts. The evaluation involves computing both the mean reward as given by the true reward function and the mean sequence-level KL\footnote{Computed as the cumulative sum of per-timestep KL divergences.} relative to the reference policy $\text{KL}\left(\pi\mid \mid \pi_\text{ref}\right)$. Our findings indicate that DPO yields the most favorable frontier—DPO attains superior rewards at lower KL values than competing methods. This result stands out for several reasons. Not only do DPO and PPO pursue the same formal objective, but DPO delivers greater efficiency and its reward/KL trade-off profile consistently surpasses that of PPO. Furthermore, DPO achieves a stronger frontier even in the case where PPO has access to exact reward information (PPO-GT).

\subsection{Can DPO scale to real preference datasets?}
\label{sec:dpo-real-datasets}
We proceed to assess the efficacy of DPO fine-tuning on tasks such as summarization and single-turn dialogue. It is well established that standard automatic metrics for summarization, like ROUGE, often do not align well with human judgments~\citep{stiennon2022learning}. Previous studies have demonstrated improvements when using PPO for fine-tuning language models on human preference data, leading to higher quality summaries. Our experimental setup involves generating model outputs on the test portion of the TL;DR summarization dataset and measuring each method's average win rate relative to reference summaries from the test set. Samples are generated for each method across a range of temperature values from 0.0 to 1.0, and the corresponding win rates are depicted in Figure~\ref{fig:frontier-tldr-main} (right). DPO, PPO, and Preferred-FT all start from the same GPT-J SFT checkpoint\footnote{\url{https://huggingface.co/CarperAI/openai_summarize_tldr_sft}} for fine-tuning. Our results indicate that DPO achieves a win rate of about 61\% at a sampling temperature of 0.0, surpassing PPO, which attains approximately 57\% at its optimal temperature of 0.0. Furthermore, DPO attains a higher best-case win rate compared to the $N$-best baseline. It is worth noting that we did not conduct an extensive search over DPO’s $\beta$ hyperparameter, suggesting the reported performance may not reflect DPO's full capability. Additionally, DPO demonstrates greater stability across different sampling temperatures, whereas PPO’s effectiveness diminishes, sometimes reverting to that of the baseline GPT-J model at elevated temperatures. Preferred-FT provides little improvement over the initial SFT model. In Section~\ref{sec:human-judgments}, we extend this comparison with direct human evaluation: at temperature 0.25, DPO-generated outputs are preferred 58\% of the time over those from PPO at temperature 0.

For single-turn dialogue evaluation, we analyze various approaches using the subset of the Anthropic HH dataset's test split \citep{bai2022training} containing a single human-assistant interaction. In assessments with GPT-4, we use the preferred responses from the test set as ground truth to determine win rates for each method. Lacking a standard SFT model for this benchmark, our pipeline initiates with a pre-trained Pythia-2.8B. We apply Preferred-FT to fine-tune a reference model on the curated completions, thereby aligning the outputs with the model's data distribution, and subsequently perform DPO training. Comparisons are also made with the top outcome among 128 Preferred-FT completions (having observed, as detailed in Appendix Figure~\ref{fig:best-of-n}, that performance plateaus at 128 for this Best of $N$ baseline), as well as a 2-shot prompt format using the unmodified Pythia-2.8B. Results indicate that DPO matches or exceeds performance relative to the strongest temperatures explored across all baselines. Additionally, an RLHF model trained via PPO on the Anthropic HH dataset\footnote{\url{https://huggingface.co/reciprocate/ppo_hh_pythia-6B}} from a reputable repository\footnote{\url{https://github.com/CarperAI/trlx/tree/main/examples/hh}} was evaluated, but no prompt or temperature setting yielded results surpassing those of the original Pythia-2.8B. Drawing on our findings from TL;DR and noting the shared reward objective of both methods, we regard the Best of 128 approach as an approximate stand-in for PPO-level outcomes. Summarizing, DPO emerges as the only efficient method that demonstrates improvements over the preferred completions in the Anthropic HH dataset, achieving performance comparable to, or exceeding, the much more resource-intensive Best of 128 baseline. Notably, as depicted in Figure~\ref{fig:dialogue-main}, DPO rapidly approaches its optimal performance.

\subsection{Generalization to a new input distribution}

To further investigate how PPO and DPO handle shifts in input distribution, we assess the policies trained in our Reddit TL;DR summarization setup on a new dataset: news articles from the CNN/DailyMail test split \citep{nallapati-etal-2016-abstractive}. The evaluation uses the optimal sampling temperatures identified for TL;DR (0 and 0.25). Table~\ref{tab:ood} summarizes these outcomes. We measure the win rate of GPT-4 when compared to reference summaries in the CNN/DailyMail dataset, following the same GPT-4 (C) evaluation protocol as used for Reddit TL;DR, with terminology adjusted from ``forum post'' to ``news article''. Across this alternative domain, DPO demonstrates a marked advantage over the PPO policy. These findings serve as preliminary evidence that DPO approaches can generalize comparably to PPO, despite not utilizing the additional set of unlabeled Reddit TL;DR prompts accessible to PPO during training.

\subsection{Validating GPT-4 judgments with human judgments}
\label{sec:human-judgments}

We implement a human evaluation study to confirm the validity of GPT-4’s judgments using results from the TL;DR summarization experiment and two distinct GPT-4 prompts. The \textbf{GPT-4 (S)} (simple) prompt queries which summary more effectively captures the essential content of the post. Conversely, the \textbf{GPT-4 (C)} (concise) prompt also incorporates conciseness as a criterion; we examine this prompt because, under \textbf{GPT-4 (S)}, GPT-4 tends to prefer summaries that are longer and more repetitive compared to human raters. The complete text of both prompts is available in Appendix~\ref{app:prompts}. For the study, we select three approaches to represent a range of sample qualities: the top-performing method (DPO, temp. 0.25), the lowest-performing (PPO, temp. 1.0), and an intermediate one (SFT, temp. 0.25), each compared against PPO with its optimal temperature under greedy sampling. Both prompts yield agreement between GPT-4 and human raters that is on par with human-human agreement, indicating that GPT-4 serves as an adequate stand-in for human judgment (note that, due to a small number of human raters, multiple ratings were collected only for DPO and PPO-1 comparisons). Overall, the win rates generated via the \textbf{GPT-4 (C)} prompt are more aligned with those observed in human assessments, and thus this prompt is employed in the principal results reported in Section~\ref{sec:dpo-real-datasets}. Further methodological details regarding the human study—including the online interface provided to raters and a list of participants—are available in Appendix~\ref{app:human-study}.

\section{Discussion}
Preference-based learning offers an effective and scalable approach for developing proficient and aligned language models. In this work, we proposed DPO, a streamlined framework for leveraging preference data to train language models without resorting to reinforcement learning. Instead of transforming preference learning into a reinforcement learning task simply to employ traditional RL methods, DPO formulates a connection between language model policies and reward functions. This enables the language model to be optimized according to human preferences \textit{directly}, utilizing a straightforward cross-entropy loss function, entirely sidestepping reinforcement learning methodologies and without sacrificing generality. Remarkably, DPO achieves comparable or improved outcomes relative to established RLHF algorithms such as those employing PPO, all with minimal hyperparameter tuning; this considerably lowers the entry barrier for leveraging human preference data in language model training.

\textbf{Limitations \& Future Work.} Our findings open several promising avenues for further investigation. One central question is how the DPO-trained policies handle out-of-distribution generalization when contrasted with approaches that utilize explicitly learned reward functions. Preliminary experiments indicate that DPO matches PPO-based approaches in this regard, yet comprehensive analyses remain necessary. For instance, further research is needed to determine if self-labeling with the DPO policy can efficiently exploit unlabeled prompt datasets. Additionally, a better understanding is needed regarding how over-optimization of the reward signal presents itself within the DPO paradigm, and whether the slight performance drop observed in Figure~\ref{fig:dialogue-main}-right is indicative of this phenomenon. Our current assessments are limited to models with up to 6B parameters; scaling DPO to much larger, cutting-edge models is an especially intriguing prospect for future work. With respect to evaluation, we observe that GPT-4’s win rates are prompt-sensitive, suggesting a need for future investigations into optimizing automated evaluation quality. Lastly, DPO’s potential extends well beyond aligning language models to human preferences, with promising applications in training generative models across various data modalities.